% !TeX root = ../tfg.tex
% !TeX encoding = utf8
%
%*******************************************************
% Summary
%*******************************************************

\section*{Resumen}
\addcontentsline{toc}{section}{Resumen}

\begin{center}
\emph{Metodologías Multivariantes para la Identificación de Patrones Biológicos}
\end{center}
\vspace{1cm}

El desarrollo de las denominadas ciencias ómicas, ha transformado la forma en que se estudian los sistemas biológicos, permitiendo analizar 
a nivel molecular grandes volúmenes de datos sobre genes, proteínas o metabolitos. En particular, la transcriptómica, centrada en el estudio 
de los ARN mensajeros presentes en una célula o tejido, se ha convertido en una herramienta esencial para entender procesos biológicos y estados 
celulares. Sin embargo, el volumen, complejidad y alta dimensionalidad de los datos generados por técnicas de secuenciación de ARN como RNA-seq 
plantean numerosos desafíos desde el punto de vista estadístico y computacional. \newline

Frente a esta situación, las técnicas multivariantes son muy adecuadas para gestionar y explorar este tipo de información. Técnicas como el 
Análisis de Componentes Principales (PCA), el Análisis Factorial, el Análisis Discriminante o el Análisis Clúster, permiten reducir la dimensionalidad, 
identificar patrones ocultos y clasificar observaciones en función de su similitud. En este trabajo se ha realizado una revisión detallada de estas 
metodologías, haciendo especial hincapié en el análisis cluster por su capacidad para detectar agrupaciones naturales en datos sin etiquetar 
(aprendizaje automático no supervisado). \newline

Datos que provienen de bases de datos como la del proyecto GEO, una de las principales fuentes públicas de datos RNA-seq . Aprovechando las capacidades 
del entorno Bioconductor y de R para el análisis multivariante, se identificó la necesidad de desarrollar una aplicación web sencilla que permitiera 
aplicar técnicas de análisis multivariante como las mencionadas anteriormente a una matriz de expresión génica. Frente a esta necesidad, se ha desarrollado 
una aplicación software en Shiny R siguiendo para ello todas las fases del ciclo de vida del software e incluyendo un manual de usuario exhaustivo con 
todas las indicaciones para utilizar la herramienta de manera efectiva por cualquier usuario. \newline

Con todo ello, se ha logrado crear un primer prototipo de solución integral que no solo facilita el análisis de datos transcriptómicos, sino que también 
mejora la accesibilidad y la usabilidad de estos métodos estadísticos en el contexto de la bioinformática.


\noindent

\newpage

\selectlanguage{english}
\begin{center}
    \emph{Multivariate Methodologies for the Identification of Biological Patterns}
    \end{center}
    \vspace{1cm}

    The development of the so-called omics sciences has transformed the way biological systems are studied, enabling the molecular-level analysis of large volumes 
    of data on genes, proteins, or metabolites. In particular, transcriptomics, focused on the study of messenger RNAs present in a cell or tissue, has become an 
    essential tool for understanding biological processes and cellular states. However, the volume, complexity, and high dimensionality of the data generated by RNA 
    sequencing techniques, such as RNA-seq, pose numerous challenges from both a statistical and computational perspective. \newline

    In response to this situation, multivariate techniques are highly suitable for managing and exploring this type of information. Techniques such as Principal 
    Component Analysis (PCA), Factor Analysis, Discriminant Analysis, and Cluster Analysis enable dimensionality reduction, the identification of hidden patterns, 
    and the classification of observations based on their similarity. This work provides a detailed review of these methodologies, with a special emphasis on cluster 
    analysis due to its ability to detect natural groupings in unlabeled data (unsupervised machine learning). \newline

    Data from databases such as the GEO project, one of the main public sources of RNA-seq data, have been key to addressing these challenges. By leveraging the 
    capabilities of the Bioconductor environment and R for multivariate analysis, the need to develop a simple web application that would allow the application of 
    multivariate analysis techniques, as mentioned earlier, to gene expression matrices was identified. In response to this need, a software application was developed 
    in Shiny R, following all the phases of the software development life cycle, and including a comprehensive user manual with detailed instructions for effectively 
    using the tool. \newline

    With this, the first prototype of an integrated solution has been created, which not only facilitates transcriptomic data analysis but also improves the 
    accessibility and usability of these statistical methods in the context of bioinformatics.
    
\selectlanguage{spanish} 
\endinput
